\par\addvspace{12pt}\Needspace{9\baselineskip}

\section{条件概率}

% ITEM: probability-example-010
\solutionitem{probability-example-010}{例6.10}
（2023年全国甲卷理科，题面整理）
某地中学生中，60\%爱好滑冰，50\%爱好滑雪，70\%至少爱好其中一项。随机调查一名学生，已知他爱好滑雪，求他也爱好滑冰的概率。
\solutionanswer
记滑冰、滑雪为\(A,B\)，则\(P(A\cap B)=0.6+0.5-0.7=0.4\)。已知滑雪后分母应取\(P(B)=0.5\)，所以\(P(A\mid B)=0.8\)。


% ITEM: probability-example-011
\solutionitem{probability-example-011}{例6.11}
（2022年天津卷，题面整理）
从含4张A的52张扑克牌中依次随机取两张，不放回。求两张都是A的概率，以及第一张是A时第二张也是A的条件概率。
\solutionanswer
第一张为A的概率\(4/52\)，在此条件下第二张仍为A的概率\(3/51=1/17\)。乘法公式给出两张均为A的概率\((4/52)(3/51)=1/221\)。


% ITEM: probability-example-012
\solutionitem{probability-example-012}{例6.12}
（据2025年北京卷改编）
甲、乙两校各随机调查100名学生，分别有80人、75人答对某四选一选择题。用频率估计答对概率。甲校学生掌握知识点时必答对，乙校学生掌握时答对概率为\(0.85\)；两校未掌握者均等可能地猜选四个选项。分别估计两校学生掌握知识点的比例，并比较大小。
\solutionanswer
用答对频率分别估计为\(0.80,0.75\)。设两校掌握比例为\(p,q\)，全概率公式给出\(0.80=p+(1-p)/4\)、\(0.75=0.85q+(1-q)/4\)。解得\(p=11/15\)、\(q=5/6\)。乙校掌握比例估计更高；不能只比较表面的答对比例，因为掌握者的答对机制也不同。


% ITEM: probability-exercise-04-01
\solutionitem{probability-exercise-04-01}{习题6.4（1）}
某盒中有4个红球、3个白球，依次不放回抽两个。已知第二个是白球，求第一个是红球的概率。
\solutionanswer
设\(R_1,W_2\)分别表示第一红、第二白。\(P(R_1\cap W_2)=(4/7)(3/6)=2/7\)，而第二球白的概率由对称性为\(3/7\)，所以条件概率为\(2/3\)。也可在第二球白后，直接考察其余六球中第一球的颜色，有四红二白。


% ITEM: probability-exercise-04-02
\solutionitem{probability-exercise-04-02}{习题6.4（2）}
本节模型：甲、乙来源比例为\(0.4,0.6\)，相应不合格率为\(0.02,0.05\)。
在本节生产线模型中，已知抽到的产品合格，求它来自乙线的概率。
\solutionanswer
合格概率\(0.4\times0.98+0.6\times0.95=0.962\)，乙线且合格的概率\(0.57\)。故\(P(\text{乙}\mid\text{合格})=0.57/0.962=285/481\)。


% ITEM: probability-exercise-04-03
\solutionitem{probability-exercise-04-03}{习题6.4（3）}
盒甲有2个红球、1个白球，盒乙有1个红球、3个白球。等可能地选一盒，再从中随机取一球。已知取出红球，求所选为盒甲的概率。若改为从全部7个球中等可能抽一个，条件概率如何改变？
\solutionanswer
等可能选盒时，甲盒红的联合概率为\((1/2)(2/3)\)，总红概率\((1/2)(2/3)+(1/2)(1/4)\)，相除得\(8/11\)。改为从七球均匀抽取时，三个红球等可能，其中两个来自甲盒，条件概率\(2/3\)。两种机制对盒子大小赋予的权重不同。


% ITEM: probability-exercise-04-04
\solutionitem{probability-exercise-04-04}{习题6.4（4）}
设\(P(B)=0.3\)、\(P(A\mid B)=0.8\)、\(P(A\mid B^c)=0.2\)。求\(P(A)\)、\(P(B\mid A)\)及\(P(B\mid A^c)\)。
\solutionanswer
\(P(A)=0.3\times0.8+0.7\times0.2=0.38\)。故\(P(B\mid A)=0.24/0.38=12/19\)，\(P(B\mid A^c)=0.3\times0.2/0.62=3/31\)。


% ITEM: probability-exercise-04-05
\solutionitem{probability-exercise-04-05}{习题6.4（5）}
（据2024年天津卷改编）从包含“整地做畦”和“田间灌溉”的5个劳动技能项目中等可能选3项。求选中前一项的概率，以及已选中前一项时也选中后一项的概率。
\solutionanswer
第一项被选中的概率\(\binom42/\binom53=3/5\)。已选第一项后，另两项从剩下四项中均匀选取，也选第二项的概率\(\binom31/\binom42=1/2\)。


% ITEM: probability-exercise-05-01
\solutionitem{probability-exercise-05-01}{习题6.5 *}
掷两枚骰子，设36个有序点数对等可能。分别考虑记录“点数和”“第一枚点数”“两个点数的有序对”所形成的分割。比较其粗细关系，并求已知点数和为8时第一枚大于3的概率。
\solutionanswer
有序对分割比另外两种都细；“点数和”与“第一枚点数”互不比对方更细：同一和可以有不同第一点数，同一第一点数也可以有不同和。同时知道两项记录，则能恢复有序对。和为8的有序对为\((2,6),(3,5),(4,4),(5,3),(6,2)\)，条件下仍等可能，第一枚大于3的有三个，故概率\(3/5\)。
\par
